% Options for packages loaded elsewhere
\PassOptionsToPackage{unicode}{hyperref}
\PassOptionsToPackage{hyphens}{url}
\PassOptionsToPackage{dvipsnames,svgnames,x11names}{xcolor}
\documentclass[
  10pt,
  fontset=fandol]{ctexart}
\usepackage{xcolor}
\usepackage[margin=16mm]{geometry}
\usepackage{amsmath,amssymb}
\setcounter{secnumdepth}{-\maxdimen} % remove section numbering
\usepackage{iftex}
\ifPDFTeX
  \usepackage[T1]{fontenc}
  \usepackage[utf8]{inputenc}
  \usepackage{textcomp} % provide euro and other symbols
\else % if luatex or xetex
  \usepackage{unicode-math} % this also loads fontspec
  \defaultfontfeatures{Scale=MatchLowercase}
  \defaultfontfeatures[\rmfamily]{Ligatures=TeX,Scale=1}
\fi
\usepackage{lmodern}
\ifPDFTeX\else
  % xetex/luatex font selection
\fi
% Use upquote if available, for straight quotes in verbatim environments
\IfFileExists{upquote.sty}{\usepackage{upquote}}{}
\IfFileExists{microtype.sty}{% use microtype if available
  \usepackage[]{microtype}
  \UseMicrotypeSet[protrusion]{basicmath} % disable protrusion for tt fonts
}{}
\makeatletter
\@ifundefined{KOMAClassName}{% if non-KOMA class
  \IfFileExists{parskip.sty}{%
    \usepackage{parskip}
  }{% else
    \setlength{\parindent}{0pt}
    \setlength{\parskip}{6pt plus 2pt minus 1pt}}
}{% if KOMA class
  \KOMAoptions{parskip=half}}
\makeatother
\setlength{\emergencystretch}{3em} % prevent overfull lines
\providecommand{\tightlist}{%
  \setlength{\itemsep}{0pt}\setlength{\parskip}{0pt}}
\usepackage{amsmath,amssymb,mathtools,needspace}
\xeCJKDeclareCharClass{CJK}{"2032,"0400->"04FF,"2160->"216B,"2460->"2473,"25A0,"25C7,"25CB,"25CF}
\IfFontExistsTF{Arial Unicode MS}{\xeCJKsetup{AutoFallBack=true}\setCJKfallbackfamilyfont{\CJKrmdefault}{Arial Unicode MS}}{}
\setcounter{secnumdepth}{0}
\setlength{\emergencystretch}{2em}
\allowdisplaybreaks
\usepackage{bookmark}
\IfFileExists{xurl.sty}{\usepackage{xurl}}{} % add URL line breaks if available
\urlstyle{same}
\hypersetup{
  pdftitle={引言},
  colorlinks=true,
  linkcolor={Maroon},
  filecolor={Maroon},
  citecolor={Blue},
  urlcolor={Blue},
  pdfcreator={LaTeX via pandoc}}

\title{引言}
\author{}
\date{}

\begin{document}
\maketitle

\subsection{8.1　引言}\label{ux5f15ux8a00}

考虑线性方程组

\[
Ax=b,
\tag{8.1.1}
\]

其中 \(A\) 为非奇异矩阵。当 \(A\) 为低阶稠密矩阵时，第 7 章所讨论的选主元素消去法是解式（8.1.1）的有效方法。但是，对于由工程技术中产生的大型稀疏矩阵方程组（\(A\) 阶数 \(n\) 很大，但零元素较多，例如 \(n\geqslant10^4\)，由某些偏微分方程数值解所产生的线性方程组），利用迭代法求解式（8.1.1）是合适的。在计算机内存和运算两方面，迭代法通常都可利用 \(A\) 中有大量零元素的特点。

本章将介绍迭代法的一些基本理论及 Jacobi 迭代法、Gauss-Seidel 迭代法、超松弛迭代法。超松弛迭代法应用很广泛。

下面举简例，以便了解迭代法的思想。

\textbf{例 8.1}　求解方程组

\[
\begin{cases}
8x_1-3x_2+2x_3=20,\\
4x_1+11x_2-x_3=33,\\
6x_1+3x_2+12x_3=36,
\end{cases}
\tag{8.1.2}
\]

记作

\[
Ax=b,
\]

其中

\[
A=\begin{bmatrix}8&-3&2\\4&11&-1\\6&3&12\end{bmatrix},\qquad
x=\begin{bmatrix}x_1\\x_2\\x_3\end{bmatrix},\qquad
b=\begin{bmatrix}20\\33\\36\end{bmatrix}.
\]

方程组的精确解是

\[
x^*=(3,2,1)^{\mathrm T}.
\]

现将式（8.1.2）改写成

\[
\begin{cases}
x_1=\dfrac18(3x_2-2x_3+20),\\[4pt]
x_2=\dfrac1{11}(-4x_1+x_3+33),\\[4pt]
x_3=\dfrac1{12}(-6x_1-3x_2+36)
\end{cases}
\tag{8.1.3}
\]

\href{../../source-images/pdf-216.jpeg}{原页图像}

或

\[
x=B_0x+f,
\]

其中

\[
B_0=\begin{bmatrix}
0&\dfrac38&-\dfrac28\\[4pt]
-\dfrac4{11}&0&\dfrac1{11}\\[4pt]
-\dfrac6{12}&-\dfrac3{12}&0
\end{bmatrix}=I-D^{-1}A,\qquad
f=\begin{bmatrix}\dfrac{20}8\\[4pt]\dfrac{33}{11}\\[4pt]\dfrac{36}{12}\end{bmatrix}=D^{-1}b.
\]

任取初始值，例如取 \(x^{(0)}=(0,0,0)^{\mathrm T}\)。将这些值代入式（8.1.3）右端（若式（8.1.3）为等式即求得方程组的解，但一般不满足），得到新的值

\[
x^{(1)}=(x_1^{(1)},x_2^{(1)},x_3^{(1)})^{\mathrm T}=(2.5,3,3)^{\mathrm T},
\]

再将 \(x^{(1)}\) 分量代入式（8.1.3）右端，得到 \(x^{(2)}\)。反复利用这个计算程序，得到一向量序列和一般的计算公式（迭代公式）

\[
x^{(0)}=\begin{bmatrix}x_1^{(0)}\\x_2^{(0)}\\x_3^{(0)}\end{bmatrix},\quad
x^{(1)}=\begin{bmatrix}x_1^{(1)}\\x_2^{(1)}\\x_3^{(1)}\end{bmatrix},\quad\cdots,\quad
x^{(k)}=\begin{bmatrix}x_1^{(k)}\\x_2^{(k)}\\x_3^{(k)}\end{bmatrix},\quad\cdots,
\]

\[
\begin{cases}
x_1^{(k+1)}=(3x_2^{(k)}-2x_3^{(k)}+20)/8,\\
x_2^{(k+1)}=(-4x_1^{(k)}+x_3^{(k)}+33)/11,\\
x_3^{(k+1)}=(-6x_1^{(k)}-3x_2^{(k)}+36)/12,
\end{cases}
\tag{8.1.4}
\]

简写为

\[
x^{(k+1)}=B_0x^{(k)}+f,
\]

其中 \(k\)（\(k=0,1,2,\cdots\)）表示迭代次数。

迭代到第 10 次有

\[
x^{(10)}=(3.000\,032,\ 1.999\,838,\ 0.999\,881\,3)^{\mathrm T},
\]

\[
\|\varepsilon^{(10)}\|_\infty=0.000\,187\quad(\varepsilon^{(10)}=x^{(10)}-x^*).
\]

从此例看出，由迭代法作出的向量序列 \(x^{(k)}\) 逐步逼近方程组的精确解 \(x^*\)。

对于任何一个方程组 \(x=Bx+f\)（由 \(Ax=b\) 变形得到的等价方程组），按迭代法作出的向量序列 \(x^{(k)}\) 是否一定逐步逼近方程组的解 \(x^*\) 呢？回答是不一定。请读者考虑，用迭代法解下述方程组

\[
\begin{cases}
x_1=2x_2+5,\\
x_2=3x_1+5.
\end{cases}
\]

对于给定方程组 \(x=Bx+f\)，设有唯一解 \(x^*\)，则

\[
x^*=Bx^*+f.
\tag{8.1.5}
\]

又设 \(x^{(0)}\) 为任取的初始向量，按下述公式构造向量序列

\[
\begin{cases}
x^{(1)}=Bx^{(0)}+f,\\
x^{(2)}=Bx^{(1)}+f,\\
\qquad\vdots\\
x^{(k+1)}=Bx^{(k)}+f,
\end{cases}
\tag{8.1.6}
\]

\begin{quote}
校注（agent 补充）：原页将 \(x^{(10)}\) 的第二分量印为 \(1.999\,838\)，误差范数印为 \(0.000\,187\)，均按原样保留。由式（8.1.4）从 \(x^{(0)}=0\) 作 10 次有理数迭代，所得向量约为 \((3.0000318141,1.9998740186,0.9998812605)^{\mathrm T}\)，实际误差的无穷范数约为 \(0.0001259814\)；即使按本页印刷的三个分量计算，误差无穷范数也应为 \(0.000162\)。第二分量及误差范数疑为原书数值错误，并非 OCR 错误。
\end{quote}

\href{../../source-images/pdf-217.jpeg}{原页图像}

其中 \(k\) 表示迭代次数。

\textbf{定义 8.1}　\(1^\circ\) 对于给定的方程组 \(x=Bx+f\)，用式（8.1.6）逐步代入求近似解的方法称为\textbf{迭代法}（或称为一阶定常迭代法，这里 \(B\) 与 \(k\) 无关）。

\(2^\circ\) 如果 \(\lim\limits_{k\to\infty}x^{(k)}\) 存在（记作 \(x^*\)），称此迭代法\textbf{收敛}，显然 \(x^*\) 就是方程组的解，否则称此迭代法\textbf{发散}。

由上述讨论，需要研究 \(\{x^{(k)}\}\) 的收敛性。引进误差向量

\[
\varepsilon^{(k+1)}=x^{(k+1)}-x^*,
\]

由式（8.1.6）减去式（8.1.5），得

\[
\varepsilon^{(k+1)}=B\varepsilon^{(k)}\quad(k=0,1,2,\cdots),
\]

递推得

\[
\varepsilon^{(k)}=B\varepsilon^{(k-1)}=\cdots=B^k\varepsilon^{(0)}.
\]

要考察 \(\{x^{(k)}\}\) 的收敛性，就要研究 \(B\) 在什么条件下有 \(\varepsilon^{(k)}\to0\)（\(k\to\infty\)），亦即要研究 \(B\) 满足什么条件时有 \(B^k\to O\)（零矩阵）（\(k\to\infty\)）。

\begin{center}\rule{0.5\linewidth}{0.5pt}\end{center}

\subsection{相关原页校注（agent 补充）}\label{ux76f8ux5173ux539fux9875ux6821ux6ce8agent-ux8865ux5145}

以下校注来自本节覆盖的原页，可能也涉及同页相邻小节；原文照录与校正说明分开保留。

\subsubsection{原书 PDF 页 217 的校注}\label{ux539fux4e66-pdf-ux9875-217-ux7684ux6821ux6ce8}

\href{../pages/pdf-217.md}{完整原页转录} · \href{../../source-images/pdf-217.jpeg}{原页图像}

校注记录：ch08-E013。

\begin{quote}
校注（agent 补充，ch08-E013）：本页式（8.2.1）下方原印 \(a_{ij}\neq0\)（\(i=1,2,\cdots,n\)），已按原样保留。这里没有给出 \(j\) 的范围；后续 Jacobi 分裂以对角元 \(a_{ii}\) 作除数，通常条件应为 \(a_{ii}\neq0\)。此处疑为原书下标错误；独立交叉复核指出后已回原图放大确认。
\end{quote}

\subsection{本节覆盖原页的校注记录}\label{ux672cux8282ux8986ux76d6ux539fux9875ux7684ux6821ux6ce8ux8bb0ux5f55}

下列记录按原页关联，含同页相邻小节的校注；计算时同时读取上文校注和完整原页。

\begin{itemize}
\tightlist
\item
  \texttt{ch08-E013}：原书 PDF 页 217
\item
  \texttt{ch08-E001}：原书 PDF 页 216
\end{itemize}

\end{document}
